% !TEX root = ../main.tex
\begin{ouverture}[cinetique]{4}{Cinétique chimique}{Partie 1 · Constitution et transformations de la matière}
{Suivi temporel\enspace·\enspace Vitesses volumiques\enspace·\enspace Facteurs cinétiques et catalyse\enspace·\enspace Loi d'ordre 1}
\objectif{Justifier le choix d'une méthode de suivi temporel d'une transformation chimique.}
\objectif{Déterminer une vitesse volumique de disparition ou d'apparition et un temps de demi-réaction.}
\objectif{Identifier des facteurs cinétiques et un catalyseur à partir de données expérimentales.}
\objectif{Reconnaître une évolution qui suit une loi de vitesse d'ordre 1.}
\objectif{Tester une loi de vitesse à l'aide d'un programme Python.}
\prerequis{Tableau d'avancement, réactif limitant, avancement final.}
\prerequis{Lois de Beer-Lambert et de Kohlrausch, gaz parfait (chapitre 2).}
\prerequis{Dérivée, fonction exponentielle et logarithme népérien.}
\end{ouverture}
\essentiel
\essSection{1}{Transformations lentes ou rapides et suivi temporel}
\begin{definition}[Transformation lente, transformation rapide]
Une transformation est \textbf{rapide} lorsqu'elle semble achevée dès que les réactifs entrent en contact : son
évolution ne peut être suivie ni à l'œil, ni avec les appareils de mesure usuels. Elle est \textbf{lente} lorsque son
évolution dure de quelques secondes à plusieurs heures, voire davantage : on peut alors mesurer la composition du
système au cours du temps. La \textbf{cinétique chimique} étudie cette évolution temporelle.
\end{definition}
\begin{propriete}[Méthodes de suivi temporel]
On suit l'évolution d'un système en mesurant, à différentes dates, une grandeur liée à la concentration d'une espèce.
\par\smallskip
{\small\renewcommand{\arraystretch}{1.2}
\begin{tabularx}{\linewidth}{@{}llY@{}}\toprule
\otTete{Méthode} & \otTete{Grandeur mesurée} & \otTete{Condition d'utilisation}\\\midrule
Spectrophotométrie & absorbance $A=k\,c$ & une espèce colorée, seule à absorber à la longueur d'onde choisie\\
Conductimétrie & conductivité $\sigma$ & des ions apparaissent ou disparaissent\\
Pression ou volume de gaz & $p$ ou $V$ & un gaz se forme ou est consommé\\
pH-métrie & pH & les ions oxonium sont formés ou consommés\\
Titrage de prélèvements & quantité titrée & chaque prélèvement subit une trempe avant le titrage\\\bottomrule
\end{tabularx}}
\end{propriete}
\begin{savoirfaire}[Choisir une méthode de suivi]
Repérer dans l'équation une espèce dont la concentration varie et qui modifie une grandeur mesurable : couleur
(spectrophotométrie), ions (conductimétrie), gaz (pression). Les méthodes physiques fournissent des mesures en
continu, sans prélèvement et sans perturber le système ; une méthode chimique impose d'arrêter la transformation dans
chaque prélèvement par une \textbf{trempe}, avant de le titrer.
\end{savoirfaire}
\essSection{2}{Vitesses volumiques et temps de demi-réaction}
\begin{definition}[Vitesses volumiques]
Pour une transformation à volume constant, la vitesse volumique de disparition d'un réactif R et la vitesse
volumique d'apparition d'un produit P sont, à la date $t$ :
\[v_\mathrm{disp}(\mathrm{R})=-\frac{\mathrm{d}[\mathrm{R}]}{\mathrm{d}t}\qquad\text{et}\qquad
v_\mathrm{app}(\mathrm{P})=\frac{\mathrm{d}[\mathrm{P}]}{\mathrm{d}t}.\]
Ces vitesses sont positives ; elles s'expriment en \unit{mol.L^{-1}.s^{-1}} ou en \unit{mol.L^{-1}.min^{-1}}.
\end{definition}
\begin{propriete}[Vitesses et nombres stœchiométriques]
Pour la réaction $\alpha\,\mathrm{A}+\beta\,\mathrm{B}\longrightarrow\gamma\,\mathrm{C}+\delta\,\mathrm{D}$,
d'avancement $x$, dans un volume $V$ constant, $[\mathrm{A}]=[\mathrm{A}]_0-\frac{\alpha\,x}{V}$, donc
\[\frac{v_\mathrm{disp}(\mathrm{A})}{\alpha}=\frac{v_\mathrm{disp}(\mathrm{B})}{\beta}=
\frac{v_\mathrm{app}(\mathrm{C})}{\gamma}=\frac{v_\mathrm{app}(\mathrm{D})}{\delta}=\frac{1}{V}\,\frac{\mathrm{d}x}{\mathrm{d}t}.\]
\end{propriete}
\begin{savoirfaire}[Déterminer une vitesse volumique]
Tracer la tangente à la courbe $[\mathrm{R}]=f(t)$ au point d'abscisse $t$ : la vitesse volumique de disparition est
l'opposée de son coefficient directeur. Pour un produit, la vitesse d'apparition est égale au coefficient directeur
de la tangente à la courbe $[\mathrm{P}]=g(t)$. En général, ces vitesses diminuent au cours du temps.
\end{savoirfaire}
\begin{definition}[Temps de demi-réaction]
Le temps de demi-réaction $t_{1/2}$ est la durée au bout de laquelle l'avancement atteint la moitié de sa valeur
finale : $x(t_{1/2})=x_\mathrm{f}/2$. Si la transformation est totale, la concentration du réactif limitant est alors
divisée par deux.
\end{definition}
\schema{04-suivi}
\begin{attention}
Le temps de demi-réaction n'est pas la moitié de la durée de la transformation : à la date $2\,t_{1/2}$, la
transformation n'est pas terminée. Il sert à comparer des cinétiques et à choisir l'intervalle entre deux mesures ;
pour une loi d'ordre 1 (section 4), plus de 99 \% du réactif a disparu au bout de $7\,t_{1/2}$.
\end{attention}
\essSection{3}{Facteurs cinétiques et catalyse}
\begin{propriete}[Facteurs cinétiques]
Un facteur cinétique est un paramètre qui modifie la vitesse d'une transformation. En général, la vitesse augmente
quand la \textbf{température} augmente et quand la \textbf{concentration des réactifs} augmente. Lorsqu'un solide
réagit, la vitesse augmente aussi avec sa surface de contact : un solide divisé réagit plus vite.
\end{propriete}
\begin{exemple}
Le réfrigérateur ralentit la dégradation des aliments ; l'autocuiseur, où l'eau bout vers 120 °C, accélère la cuisson.
Pour figer la composition d'un prélèvement avant de le titrer, on réalise une trempe : on le verse dans de l'eau
glacée, ce qui abaisse à la fois la température et les concentrations.
\end{exemple}
\begin{definition}[Catalyseur]
Un catalyseur est une espèce qui \textbf{accélère} une transformation sans être consommée : il est régénéré et
n'apparaît pas dans l'équation de la réaction. Il ne modifie pas l'état final du système. La catalyse est
\textbf{homogène} si le catalyseur et les réactifs forment une seule phase, \textbf{hétérogène} dans le cas contraire
(souvent un solide, comme le platine), et \textbf{enzymatique} si le catalyseur est une enzyme, protéine à la fois
très efficace et très sélective.
\end{definition}
\schema{04-facteurs}
\essSection{4}{Loi de vitesse d'ordre 1}
\begin{definition}[Loi de vitesse d'ordre 1]
La disparition d'un réactif R suit une loi de vitesse d'ordre 1 si, à chaque instant,
\[v_\mathrm{disp}(\mathrm{R})=k\,[\mathrm{R}],\]
où $k$ est la constante de vitesse, positive, exprimée en \unit{s^{-1}} ou en \unit{min^{-1}} ; elle dépend de la température.
\end{definition}
\begin{theoreme}[Évolution de la concentration]
Si la loi de vitesse est d'ordre 1, alors
\[[\mathrm{R}](t)=[\mathrm{R}]_0\,\mathrm{e}^{-kt},\qquad\ln\frac{[\mathrm{R}]}{[\mathrm{R}]_0}=-k\,t
\qquad\text{et}\qquad t_{1/2}=\frac{\ln2}{k}.\]
Le temps de demi-réaction ne dépend pas de la concentration initiale : la concentration est divisée par deux à
chaque nouvelle durée $t_{1/2}$.
\end{theoreme}
\begin{demonstration}
Comme $v_\mathrm{disp}(\mathrm{R})=-\frac{\mathrm{d}[\mathrm{R}]}{\mathrm{d}t}$, la fonction $t\mapsto[\mathrm{R}](t)$
vérifie l'équation différentielle $y'=-k\,y$, dont les solutions sont les fonctions $t\mapsto C\,\mathrm{e}^{-kt}$.
La condition initiale impose $C=[\mathrm{R}]_0$. À la date $t_{1/2}$, $[\mathrm{R}]=[\mathrm{R}]_0/2$, d'où
$\mathrm{e}^{-k\,t_{1/2}}=\frac12$, puis $k\,t_{1/2}=\ln2$.
\end{demonstration}
\schema{04-ordre1}
\begin{savoirfaire}[Reconnaître une loi de vitesse d'ordre 1]
À partir de mesures, vérifier l'une des propriétés suivantes : la vitesse volumique de disparition est proportionnelle
à $[\mathrm{R}]$ (les points sont alignés avec l'origine) ; $\ln([\mathrm{R}]/[\mathrm{R}]_0)$ est une fonction
affine du temps ; les temps de demi-réaction successifs sont égaux et ne dépendent pas de la concentration initiale.
\end{savoirfaire}
\begin{codecote}[0.58]
\begin{pythoncol}
import matplotlib.pyplot as plt
t = [0, 5, 10, 15, 20, 25, 30, 35, 40]  # min
c = [10.0, 8.71, 7.58, 6.60, 5.74,
     5.00, 4.35, 3.79, 3.30]  # mmol/L
V, C = [], []
for i in range(1, len(t) - 1):
    v = -(c[i+1] - c[i-1]) / (t[i+1] - t[i-1])
    V.append(v)
    C.append(c[i])
plt.plot(C, V, "o")
plt.xlabel("[R] (mmol/L)")
plt.ylabel("v (mmol/L/min)")
plt.show()
\end{pythoncol}
\cote
\fontsize{9.5}{12.5}\selectfont
\capacite{Capacité numérique}\par\smallskip
La vitesse à la date $t_i$ est approchée par le taux de variation de la concentration entre les dates voisines
$t_{i-1}$ et $t_{i+1}$ : la boucle commence donc au deuxième point et s'arrête à l'avant-dernier.
Les points obtenus sont alignés avec l'origine : $v=k\,[\mathrm{R}]$. La loi de vitesse est d'ordre 1 et $k$ est le
coefficient directeur de la droite.
\end{codecote}
\schema{04-python-vitesse}
\begin{rediger}
Pour une vitesse volumique, préciser la date, tracer la tangente, choisir deux points éloignés sur cette droite et
donner le résultat avec son unité. Pour conclure à une loi d'ordre 1, citer le critère vérifié (« les points sont
alignés sur une droite passant par l'origine ») plutôt que « la courbe ressemble à une exponentielle ».
\end{rediger}
\exercices

\rubrique{Découverte}
\begin{blocLivre}{ex-4-1}
\begin{exercice}[1]{Lente ou rapide ?}
On observe les transformations suivantes.
\begin{questions}[label=\otLabel{\Alph*.},leftmargin=1.9em,itemsep=0pt,parsep=0.15em]
\item Un précipité blanc de chlorure d'argent apparaît dès que l'on verse du nitrate d'argent dans l'eau salée.
\item Un clou en fer plongé dans de l'eau salée se couvre de rouille en quelques jours.
\item Une solution violette de permanganate de potassium, mélangée à de l'acide oxalique, se décolore en quelques minutes.
\item Le pH d'une solution d'acide chlorhydrique varie dès que l'on y verse de la soude.
\item Un ruban de magnésium plongé dans l'acide chlorhydrique disparaît en une minute, avec un dégagement de dihydrogène.
\end{questions}
\begin{questions}
\item Définir une transformation lente et une transformation rapide, puis classer les transformations A à E.
\item Pour les transformations C et E, proposer une grandeur physique à mesurer au cours du temps pour suivre
l'évolution du système.
\item Pourquoi la réaction support d'un titrage doit-elle être rapide ?
\end{questions}
\end{exercice}
\end{blocLivre}
\begin{blocLivre}{ex-4-2}
\begin{exercice}[1]{Lire un temps de demi-réaction}
On mélange une solution d'iodure de potassium (\ce{K+}, \ce{I-}), en excès, et une solution de peroxodisulfate de
sodium. La réaction \ce{S2O8^2- + 2 I- -> 2 SO4^2- + I2}, totale, est suivie par spectrophotométrie : le graphique
donne la concentration du diiode formé.
\schema{04-ex02}
\begin{questions}
\item Définir le temps de demi-réaction.
\item Déterminer la concentration finale en diiode, puis le temps de demi-réaction.
\item Les ions peroxodisulfate constituent le réactif limitant. Quelle était leur concentration initiale dans le mélange ?
\item Lire $[\ce{I2}]$ à $t=\qty{30}{min}$. Est-ce le double de sa valeur à $t=\qty{15}{min}$ ? Interpréter.
\end{questions}
\end{exercice}
\end{blocLivre}
\begin{blocLivre}{ex-4-3}
\begin{exercice}[1]{Exploiter des tangentes}
Le graphique représente la concentration d'un réactif R au cours d'une transformation. On a tracé les tangentes
T$_0$ et T$_{1/2}$ à la courbe aux dates $t=0$ et $t=t_{1/2}$.
\schema{04-ex03}
\begin{questions}
\item Définir la vitesse volumique de disparition de R. Comment la détermine-t-on graphiquement ?
\item Déterminer le temps de demi-réaction.
\item Calculer la vitesse volumique de disparition de R aux dates $t=0$ et $t=t_{1/2}$, en \unit{mol.L^{-1}.min^{-1}}.
\item Calculer le quotient $v_\mathrm{disp}(\mathrm{R})/[\mathrm{R}]$ à ces deux dates. Que constate-t-on ?
\end{questions}
\end{exercice}
\end{blocLivre}
\begin{blocLivre}{ex-4-4}
\begin{exercice}[1]{Identifier des facteurs cinétiques}
On étudie la réaction \ce{H2O2 + 2 I- + 2 H3O+ -> I2 + 4 H2O}. Dans les quatre expériences, l'eau oxygénée, de
concentration initiale \qty{1,0e-3}{mol.L^{-1}}, est le réactif limitant.
\begin{center}\small\renewcommand{\arraystretch}{1.2}
\begin{tabular}{@{}l*{4}{c}@{}}\toprule
Expérience & 1 & 2 & 3 & 4\\\midrule
$[\ce{I-}]_0$ (mol$\cdot$L$^{-1}$) & 0,10 & 0,20 & 0,10 & 0,10\\
Température (°C) & 20 & 20 & 35 & 20\\
Ajout & -- & -- & -- & ions molybdate\\
$t_{1/2}$ (min) & 12 & 6 & 4 & 1\\\bottomrule
\end{tabular}\end{center}
\begin{questions}
\item Quelles expériences faut-il comparer pour étudier l'influence de la concentration en ions iodure ? de la
température ? Conclure dans chaque cas.
\item Les ions molybdate, dissous, se retrouvent intacts en fin de réaction. Quel est leur rôle ? La catalyse est-elle
homogène ou hétérogène ?
\item La quantité finale de diiode est-elle la même dans les quatre expériences ? Justifier.
\end{questions}
\end{exercice}
\end{blocLivre}
\begin{blocLivre}{ex-4-5}
\begin{exercice}[1]{Un catalyseur solide}
On verse 50 mL d'eau oxygénée dans deux erlenmeyers. Dans le premier, on ajoute une pointe de spatule de dioxyde de
manganèse \ce{MnO2} en poudre, solide noir : un abondant dégagement de dioxygène se produit aussitôt. Dans le second,
rien n'est observé. En fin de réaction, le solide noir est filtré, séché et pesé : sa masse n'a pas varié.
\begin{questions}
\item Écrire l'équation de la décomposition du peroxyde d'hydrogène \ce{H2O2} en eau et en dioxygène.
\item Quels arguments montrent que le dioxyde de manganèse est un catalyseur de cette réaction ?
\item La catalyse est-elle homogène ou hétérogène ?
\item Le dégagement est beaucoup plus lent avec un seul gros grain de \ce{MnO2} de même masse. Proposer une explication.
\item Le volume total de dioxygène recueilli dépend-il de la forme sous laquelle on introduit le catalyseur ?
\end{questions}
\end{exercice}
\end{blocLivre}
\rubrique{Entraînement}
\begin{blocLivre}{ex-4-6}
\begin{exercice}[2]{Suivi par spectrophotométrie}
À la date $t=0$, on mélange $V_1=\qty{10,0}{mL}$ d'eau oxygénée de concentration $c_1=\qty{1,0e-2}{mol.L^{-1}}$ et
$V_2=\qty{40,0}{mL}$ d'une solution acidifiée d'iodure de potassium de concentration $c_2=\qty{0,10}{mol.L^{-1}}$.
La réaction \ce{H2O2 + 2 I- + 2 H3O+ -> I2 + 4 H2O}, totale, forme du diiode, seule espèce colorée. Un spectrophotomètre
mesure l'absorbance du mélange, liée à la concentration du diiode par $A=k\,[\ce{I2}]$, avec $k=\qty{2,5e2}{L.mol^{-1}}$.
\begin{center}\small\renewcommand{\arraystretch}{1.2}
\begin{tabular}{@{}l*{12}{c}@{}}\toprule
$t$ (min) & 0 & 2 & 4 & 6 & 8 & 10 & 15 & 20 & 25 & 30 & 40 & 50\\
$A$ & 0 & 0,063 & 0,122 & 0,172 & 0,213 & 0,250 & 0,324 & 0,378 & 0,409 & 0,441 & 0,468 & 0,484\\\bottomrule
\end{tabular}\end{center}
\begin{questions}
\item Les ions oxonium sont en excès. Montrer que l'eau oxygénée est le réactif limitant. En déduire la concentration
finale en diiode et l'absorbance finale attendue.
\item Calculer $[\ce{I2}]$ à chaque date et tracer la courbe $[\ce{I2}]=f(t)$.
\item Déterminer le temps de demi-réaction.
\item Déterminer graphiquement la vitesse volumique d'apparition du diiode à $t=0$. Comment évolue-t-elle ensuite ? Pourquoi ?
\item Quelle relation lie la vitesse volumique de disparition des ions iodure à la vitesse volumique d'apparition du diiode ?
\end{questions}
\end{exercice}
\end{blocLivre}
\begin{blocLivre}{ex-4-7}
\begin{exercice}[2]{Choisir une méthode de suivi}
On souhaite suivre l'évolution de quatre transformations lentes, réalisées en solution aqueuse.
\begin{questions}[label=\otLabel{\Alph*.},leftmargin=1.9em]
\item \ce{(CH3)3C-Cl + 2 H2O -> (CH3)3C-OH + H3O+ + Cl-}
\item \ce{Zn(s) + 2 H3O+(aq) -> Zn^2+(aq) + H2(g) + 2 H2O(l)}, dans un flacon fermé
\item \ce{2 MnO4- + 5 H2C2O4 + 6 H3O+ -> 2 Mn^2+ + 10 CO2 + 14 H2O}
\item \ce{CH3COOC2H5 + HO- -> CH3COO- + C2H5OH}, avec de la soude (\ce{Na+}, \ce{HO-})
\end{questions}
\begin{donnees}
L'ion permanganate est la seule espèce colorée (violette) ; son spectre présente un maximum vers 530 nm.
$\lambda$ (mS$\cdot$m$^2\cdot$mol$^{-1}$) : \ce{HO-} : 19,9 ; \ce{CH3COO-} : 4,09 ; \ce{Na+} : 5,01.
\end{donnees}
\begin{questions}
\item Pour chaque transformation, proposer une méthode physique de suivi et la justifier.
\item Pour la transformation D, prévoir le sens de variation de la conductivité au cours du temps.
\item Pour la transformation C, à quelle longueur d'onde faut-il régler le spectrophotomètre ? Pourquoi ?
\item Citer deux avantages des méthodes physiques sur le titrage de prélèvements successifs.
\end{questions}
\end{exercice}
\end{blocLivre}
\begin{blocLivre}{ex-4-8}
\begin{exercice}[2]{Suivi par titrages successifs}
À la date $t=0$, on ajoute quelques gouttes d'une solution de chlorure de fer(III) à 100 mL d'eau oxygénée : le
peroxyde d'hydrogène se décompose selon la réaction \ce{2 H2O2 -> 2 H2O + O2}. À différentes dates, on prélève
10,0 mL du mélange, que l'on verse aussitôt dans 50 mL d'eau glacée, puis on titre le peroxyde d'hydrogène restant
par une solution acidifiée de permanganate de potassium de concentration $c=\qty{2,00e-2}{mol.L^{-1}}$.
\begin{center}\small\renewcommand{\arraystretch}{1.2}
\begin{tabular}{@{}l*{8}{c}@{}}\toprule
$t$ (min) & 0 & 5 & 10 & 15 & 20 & 30 & 40 & 60\\
$V_\mathrm{E}$ (mL) & 16,0 & 12,0 & 9,0 & 6,7 & 5,1 & 2,8 & 1,6 & 0,5\\\bottomrule
\end{tabular}\end{center}
\begin{donnees}
Réaction support du titrage : \ce{2 MnO4- + 5 H2O2 + 6 H3O+ -> 2 Mn^2+ + 5 O2 + 14 H2O}.
\end{donnees}
\begin{questions}
\item Pourquoi verse-t-on chaque prélèvement dans de l'eau glacée ? Quel est le rôle des ions fer(III) ?
\item Montrer que $[\ce{H2O2}]=\num{5.00e-3}\times V_\mathrm{E}$, avec $V_\mathrm{E}$ en mL et $[\ce{H2O2}]$ en
\unit{mol.L^{-1}}. Calculer $[\ce{H2O2}]$ à chaque date.
\item Tracer la courbe $[\ce{H2O2}]=f(t)$ et déterminer le temps de demi-réaction.
\item Calculer le quotient $[\ce{H2O2}](t+\qty{5}{min})/[\ce{H2O2}](t)$ pour $t=0$, 5 et 10 min. Montrer que ce
résultat est compatible avec une loi de vitesse d'ordre 1 et en déduire la constante de vitesse.
\end{questions}
\end{exercice}
\end{blocLivre}
\begin{blocLivre}{ex-4-9}
\begin{exercice}[2]{La caféine dans l'organisme}
Après l'absorption d'une tasse de café, la concentration de caféine dans le plasma sanguin est maximale au bout
d'une heure environ ; la caféine est ensuite éliminée, principalement par le foie. On mesure chez un adulte :
\begin{center}\small\renewcommand{\arraystretch}{1.2}
\begin{tabular}{@{}l*{6}{c}@{}}\toprule
$t$ (h) & 1 & 2 & 4 & 6 & 8 & 12\\
$C$ (mg$\cdot$L$^{-1}$) & 3,58 & 3,14 & 2,38 & 1,81 & 1,37 & 0,77\\\bottomrule
\end{tabular}\end{center}
\begin{docu}{1}{Une grande variabilité}
Chez l'adulte en bonne santé, la demi-vie de la caféine, c'est-à-dire la durée au bout de laquelle sa concentration
plasmatique est divisée par deux, vaut en moyenne 4 h environ, avec des écarts de 2 à 8 h selon l'âge, la masse,
la grossesse, la prise de médicaments ou l'état du foie.
\source{D'après l'Autorité européenne de sécurité des aliments (EFSA), « La caféine », 2015.}
\end{docu}
\begin{questions}
\item Calculer $\ln C$ (avec $C$ en mg$\cdot$L$^{-1}$) à chaque date et représenter $\ln C$ en fonction de $t$.
\item Montrer que l'élimination de la caféine suit une loi de vitesse d'ordre 1 et déterminer la constante de vitesse $k$.
\item En déduire la demi-vie de la caféine chez cette personne. Est-elle cohérente avec le document ?
\item Cette personne a bu son café à 17 h. Quelle part de la concentration maximale reste-t-il à 23 h ? Commenter.
\end{questions}
\end{exercice}
\end{blocLivre}
\begin{blocLivre}{ex-4-10}
\begin{exercice}[2]{Une transformation qui s'accélère}
On mélange à 20 °C une solution acidifiée de permanganate de potassium et une solution d'acide oxalique
\ce{H2C2O4} en excès. La réaction
\[\ce{2 MnO4- + 5 H2C2O4 + 6 H3O+ -> 2 Mn^2+ + 10 CO2 + 14 H2O}\]
est suivie par spectrophotométrie, ce qui donne la concentration des ions permanganate.
\schema{04-ex10}
\begin{questions}
\item Comment la vitesse volumique de disparition des ions permanganate évolue-t-elle au cours du temps ? Déterminer
sa valeur maximale.
\item Cette évolution peut-elle s'expliquer par la seule influence de la concentration des réactifs ?
\item On admet que les ions manganèse(II) formés catalysent la réaction : on parle d'autocatalyse. Interpréter
l'allure de la courbe.
\item Établir la relation entre les vitesses volumiques de disparition de \ce{MnO4-} et de \ce{H2C2O4} et la vitesse
volumique d'apparition de \ce{CO2}.
\item Quelle serait l'allure de la courbe si l'on ajoutait des ions \ce{Mn^2+} au mélange initial ?
\end{questions}
\end{exercice}
\end{blocLivre}
\begin{blocLivre}{ex-4-11}
\begin{exercice}[2]{Programmer un test de loi de vitesse}
Le cristal violet, noté \ce{CV+}, réagit avec les ions hydroxyde, en large excès, en formant une espèce incolore.
Sa concentration, déterminée par spectrophotométrie, est relevée toutes les 4 min et exploitée par le programme suivant.
\begin{python}
t = [0, 4, 8, 12, 16, 20, 24, 28, 32, 36, 40]     # en min
c = [20.00, 18.08, 16.32, 14.70, 13.46, 12.18,
     11.00, 9.94, 8.97, 8.13, 7.36]               # en µmol/L
V, C = [], []
for i in range(1, len(t) - 1):
    v = -(c[i+1] - c[i-1]) / (t[i+1] - t[i-1])
    V.append(v)
    C.append(c[i])
print([round(v, 3) for v in V])
print([round(V[i] / C[i], 4) for i in range(len(V))])
\end{python}
\begin{console}
[0.46, 0.422, 0.357, 0.315, 0.308, 0.28, 0.254, 0.226, 0.201]
[0.0254, 0.0259, 0.0243, 0.0234, 0.0252, 0.0255, 0.0255, 0.0252, 0.0248]
\end{console}
\begin{questions}
\item \capacite{Capacité numérique}Quelle grandeur la variable \code{v} calcule-t-elle de façon approchée ? Pourquoi
la boucle ne porte-t-elle ni sur le premier ni sur le dernier point ?
\item Que contient la seconde liste affichée ? Montrer que les résultats sont compatibles avec une loi de vitesse
d'ordre 1 et estimer la constante de vitesse $k$.
\item En déduire le temps de demi-réaction, puis le vérifier à l'aide des mesures.
\item Écrire les instructions qui représentent la vitesse en fonction de la concentration. Quelle allure attend-on ?
\end{questions}
\end{exercice}
\end{blocLivre}
\rubrique{Objectif bac}
\begin{blocLivre}{ex-4-12}
\begin{exercice}[3]{Décolorer un colorant alimentaire}
\begin{docu}{1}{Une décoloration à suivre}
Le colorant bleu E133 (bleu brillant FCF), présent dans de nombreux sirops et confiseries, est décoloré par l'eau de
Javel, qui contient des ions hypochlorite \ce{ClO-}. Dans deux cuves contenant la même concentration initiale de
colorant, on ajoute de l'eau de Javel en large excès, à deux concentrations différentes, puis on mesure l'absorbance
à 630 nm au cours du temps.
\end{docu}
\schema{04-ex12}
\begin{donnees}
Le colorant E133 absorbe dans le rouge orangé ($\lambda_{\max}=\qty{630}{nm}$) ; les produits de la réaction sont
incolores. Dans chaque expérience, la concentration en ions hypochlorite reste pratiquement constante.
\end{donnees}
\begin{questions}
\item Justifier le choix de la spectrophotométrie et celui de la longueur d'onde de travail.
\item Déterminer le temps de demi-réaction de chaque expérience. Quel facteur cinétique est mis en évidence ?
\item \initiative Montrer que, dans l'expérience 1, la disparition du colorant suit une loi de vitesse d'ordre 1 et
déterminer la constante de vitesse $k_1$.
\item On admet que la constante de vitesse est proportionnelle à la concentration en ions hypochlorite. Est-ce cohérent
avec les mesures ? Prévoir le temps de demi-réaction pour \mbox{$[\ce{ClO-}]=\qty{2,0e-2}{mol.L^{-1}}$}.
\end{questions}
\end{exercice}
\end{blocLivre}
\begin{blocLivre}{ex-4-13}
\begin{exercice}[3]{Catalase, ions fer(III) ou rien ?}
\begin{docu}{1}{Une enzyme protectrice}
La catalase est une enzyme présente dans presque toutes les cellules vivantes (foie, pomme de terre…). Elle protège
les cellules en décomposant le peroxyde d'hydrogène, toxique, formé au cours du métabolisme :
\ce{2 H2O2 -> 2 H2O + O2}.
\end{docu}
\begin{docu}{2}{Protocole}
Dans quatre erlenmeyers, on place 20,0 mL d'une même eau oxygénée, puis on ajoute à $t=0$ : A : 1~mL d'eau
distillée ; B : 1~mL de solution de chlorure de fer(III) ; C : 1~mL d'extrait de foie, contenant de la catalase ;
D : 1~mL du même extrait, préalablement porté à ébullition pendant 10~min. On mesure le volume de dioxygène dégagé,
à 20 °C et sous la pression atmosphérique.
\end{docu}
\schema{04-ex13}
\begin{donnees}
Volume molaire des gaz dans les conditions de l'expérience : $V_\mathrm{m}=\qty{24,0}{L.mol^{-1}}$.
\end{donnees}
\begin{questions}
\item Quel est le rôle de l'expérience A ?
\item Préciser le type de catalyse mis en jeu dans les expériences B et C. Quel catalyseur est le plus efficace ?
\item Interpréter le résultat de l'expérience D.
\item \initiative Déterminer la concentration initiale de l'eau oxygénée.
\item Quel volume de dioxygène l'expérience A finirait-elle par produire ? Justifier.
\end{questions}
\end{exercice}
\end{blocLivre}
\begin{blocLivre}{ex-4-14}
\begin{exercice}[3]{Décontaminer des lentilles de contact}
\begin{docu}{1}{Un système « peroxyde »}
Certaines solutions d'entretien des lentilles contiennent 3 \% en masse de peroxyde d'hydrogène, qui détruit
bactéries et champignons mais irrite fortement l'œil. Dans l'étui de trempage, un disque recouvert de platine catalyse
la décomposition \ce{2 H2O2 -> 2 H2O + O2}. Les lentilles peuvent être portées lorsque la teneur en peroxyde est
devenue inférieure à 0,01 \%. Le fabricant impose un trempage d'au moins 6 h.
\end{docu}
\schema{04-ex14}
\begin{docu}{2}{Modèle cinétique}
Au contact d'un disque neuf, la décomposition suit une loi de vitesse d'ordre 1, de temps de demi-réaction
$t_{1/2}=\qty{30}{min}$. Au fil des utilisations, des dépôts réduisent l'activité du disque.
\end{docu}
\begin{questions}
\item Quel est le rôle du disque de platine ? De quel type de catalyse s'agit-il ? Pourquoi le couvercle est-il muni d'un évent ?
\item \initiative Avec un disque neuf, la durée minimale de trempage imposée par le fabricant est-elle suffisante ?
\item Quelle est la teneur en peroxyde d'hydrogène au bout de 6 h ?
\item Avec un disque usagé, le temps de demi-réaction passe à 45 min. Les lentilles peuvent-elles être portées au bout
de 6 h ? Pourquoi le fabricant demande-t-il de changer d'étui à chaque nouveau flacon ?
\end{questions}
\end{exercice}
\end{blocLivre}
\begin{blocLivre}{ex-4-15}
\begin{exercice}[3]{Hydrolyse du chlorure de tertiobutyle}
\begin{docu}{1}{Une transformation qui forme des ions}
Dans un mélange d'eau et d'acétone, le 2-chloro-2-méthylpropane, ou chlorure de tertiobutyle, noté RCl, de
concentration initiale $c_0$, réagit avec l'eau selon la réaction totale
\[\ce{(CH3)3C-Cl + 2 H2O -> (CH3)3C-OH + H3O+ + Cl-}.\]
On la suit par conductimétrie à 25 °C ; au bout de plusieurs heures, la conductivité ne varie plus et vaut
$\sigma_\infty=\qty{1,62}{mS.cm^{-1}}$.
\end{docu}
\schema{04-ex15}
\begin{docu}{2}{Exploitation des mesures}
\begin{python}
import numpy as np
t = np.arange(100, 1300, 100)          # en s
s = np.array([0.337, 0.604, 0.809, 0.977, 1.119, 1.211, 1.295,
              1.365, 1.418, 1.461, 1.493, 1.523])   # en mS/cm
y = np.log(1 - s / 1.62)
a, b = np.polyfit(t, y, 1)
print(round(a, 6), round(b, 3))
\end{python}
\begin{console}
-0.002328 0.006
\end{console}
\end{docu}
\begin{questions}
\item Justifier le choix de la conductimétrie pour suivre cette transformation.
\item À l'aide de la loi de Kohlrausch, montrer que $[\mathrm{RCl}]=c_0\left(1-\sigma/\sigma_\infty\right)$.
\item \capacite{Capacité numérique}Interpréter le résultat affiché : la loi de vitesse est-elle d'ordre 1 ? Donner la
valeur de la constante de vitesse.
\item En déduire le temps de demi-réaction et le vérifier sur le graphique.
\item \initiative À 35 °C, le temps de demi-réaction vaut 120 s. Par quel facteur la constante de vitesse a-t-elle été
multipliée ? Interpréter.
\end{questions}
\end{exercice}
\end{blocLivre}
